% Problem Template
% =============================================================================
% Copy this file into:
%   latex/sections/<domain>/problems/lc-XXX-problem-name.tex
%
% Fill in the metadata below and remove any sections not required.
% =============================================================================

% STATUS: draft | reviewed | final
% PROMOTE: none | exemplars | quant-prep
% TAGS: arrays, hashing, dp, graphs, etc.
% DIFFICULTY: easy | medium | hard
% SOURCE: LeetCode XXX | Custom | Other

\ProblemTitle{<Problem Title>}
\label{prob:<domain>-<short-name>}

\ProblemSummary
<Concise problem statement in your own words. Focus on what is being computed,
not how. Avoid examples unless necessary for clarity.>

% -----------------------------------------------------------------------------
% Optional: Author Thinking (excluded from exemplar builds)
% -----------------------------------------------------------------------------
\begin{Thinking}
<Informal reasoning, failed approaches, intuition-building, or mental models.
This section is intentionally exploratory and may be omitted in promoted or
exemplar versions of the book.>
\end{Thinking}

\EdgeCases
<One paragraph explaining any edge cases or considerations that might affect a potential algorithm based on the inputs provided.>

\KeyIdea
<One paragraph explaining the core insight or invariant that makes the problem
tractable. This should read like the “aha” moment, not a step-by-step solution.>

% -----------------------------------------------------------------------------
% Algorithm
% -----------------------------------------------------------------------------
\Algorithm
\begin{CodeBlock}[style=pseudo,caption={Pseudocode}]{<line count>}
<Write clear, language-agnostic pseudocode here.
Use indentation to reflect control flow.
Avoid implementation-specific details.>
\end{CodeBlock}

% -----------------------------------------------------------------------------
% Correctness
% -----------------------------------------------------------------------------
\Correctness
<Explain why the algorithm is correct.
Reference invariants, loop structure, or structural induction as appropriate.
This should read like a proof sketch, not a narrative description.>

% -----------------------------------------------------------------------------
% Complexity
% -----------------------------------------------------------------------------
\Complexity
<State time and space complexity using asymptotic notation.
Be explicit about what n represents.>

% -----------------------------------------------------------------------------
% Implementation
% -----------------------------------------------------------------------------
\bigskip
\noindent\textbf{C++ Implementation}

\lstinputlisting[
  style=cpp,
  caption={C++ Solution: <Problem Title>}
]{../problems/<domain>/<difficulty>/lc-XXX-problem-name/solution.cpp}

\bigskip
\noindent\textbf{Python Implementation}

\lstinputlisting[
  language=Python,
  backgroundcolor=\color{codebg},
  basicstyle=\ttfamily\small,
  frame=single,
  breaklines=true,
  caption={Python Solution: <Problem Title>}
]{../problems/<domain>/<difficulty>/lc-XXX-problem-name/solution.py}

\ProblemDivider
